\section{Anticipated Objections}
\label{sec:objections}

A position paper that does not name the objections it expects has
not committed to a position. The six objections below are the
objections the author has heard, anticipates, or has had to argue
against in writing F(AI)\textsuperscript{2}R itself. Each is
restated as charitably as we know how, and each receives an answer
that the position can defend. Where the answer is conditional, we
mark the condition; where the objection lands, we concede in
\S\ref{sec:fm-residual}.

\paragraph{``This is just FAIR with extra steps.''}
The objection is that the four canonical FAIR axes already cover
what F(AI)\textsuperscript{2}R proposes, and the new artefact
classes (transcripts, prompts, ladders) are absorbed into existing
sub-principles without rewording them. The answer is empirical:
two prior re-readings (FAIR4RS~\cite{chuehong2022fair4rs},
FAIR4ML~\cite{ravi2024fair4ml}) judged that adding sub-principles
for new artefact classes was preferable to absorbing them; we
follow that precedent. The audit pass on \textbf{Reusable} is the
load-bearing case: replaying the writing pipeline is materially
different from replaying a build, and the FAIR4RS framing for
software builds does not yet capture the cost of replaying an
LLM-mediated authoring loop.

\paragraph{``Adopting eight practices on day one is too expensive.''}
The objection is one of overhead. Our answer is that eight is the
\emph{minimum}: the failure-modes table (Table~\ref{tab:failure-modes})
shows that dropping any one practice exposes a specific failure
mode that the others do not cover. The verification ladder, in
particular, looks expensive on first pass and saves the
human-reader budget on the second --- the $\textsc{ai-confirmed}$
rung was introduced exactly to attack the verification-overhead
bottleneck (\S\ref{sec:ladder}).

\paragraph{``Provenance graphs add bureaucracy without preventing
fraud.''}
The objection is that an author willing to fabricate a citation is
willing to fabricate a triple. The answer is that the graph does
not aim to prevent fraud; it aims to make fraud
\emph{detectable by inspection}. A bad-faith adoption of the pattern
leaves a graph whose edges contradict the prose, whose claims
reference sources that do not exist, or whose
\texttt{prov:wasGeneratedBy} activities have no transcript record.
Each of those is a checkable defect. The contrast is with
unstructured AI-assisted writing, where the same fraud leaves no
artefact a third party can point at.

\paragraph{``You cannot check the model's reasoning, only its
output.''}
The objection is that LLMs are stochastic and that any artefact-level
discipline is a Potemkin pass over an opaque generator. The answer
is two-part. \emph{Yes}, we cannot inspect the model's reasoning;
we never claim otherwise. \emph{No}, the discipline is not made
useless by that limit: it makes the \emph{output's relationship to
the input} checkable, even when the model's internals are not.
That is what the verification ladder, the Source Analyzer hand-off,
and the human-reader rung are for. The sustainability section
(\S\ref{sec:sustainability}) is honest about the limit.

\paragraph{``This privileges well-resourced groups.''}
The objection is real and we name it without flinching
(\S\ref{sec:sustainability}, equity paragraph). Research groups
without access to frontier models cannot apply F(AI)\textsuperscript{2}R
at the same effective rung. F(AI)\textsuperscript{2}R describes the
work but does not equalise access to it. Two partial mitigations
follow from the verification ladder itself: smaller-model fallbacks
at audit-only stages (layout, hedging-density, list-of-lists
detection) that do not need frontier-grade reasoning, and the
structural property that an honestly-declared lower-rung
distribution at submission time is preferable to a forged
higher-rung distribution from a richer group. The pattern's
disclosure rules are written so that the difference is legible to
the reviewer.

\paragraph{``A position paper should not also be a worked example.''}
The objection is generic: position papers prescribe; case-study
papers exemplify; mixing the two is bad form. The answer is
intentional. The fifth integrated practice
(\S\ref{sec:eight}, item~5) requires that a paper which prescribes
a discipline have been written under that discipline. The
recursive case study is not an aesthetic choice; it is the way the
position can be checked. The Author's Note up front is the
invitation to read it that way.

\bigskip
\noindent The objections we cannot answer --- pre-publication
leakage, audience fatigue, adversarial review --- sit in
\S\ref{sec:fm-residual} as residual.
